\gddchapter{Analysis}

\section{Thematic Analysis Framework}
Six phase model to move from theoretical raw audio to theoretical insights:
This is just a scaffold that will be filled out with actual data from this session:


% \begin{enumerate}
%     \item Familiarisation
    
%     Listen to the recorings while reading chapter 5 to notice intiial patterns

%     \item Initial Coding
    
%     Generate labels for specific UX moments such as "Command frustration", "modal confusion" etc.

%     \item Searching for Themes
    
%     Group codes into broader categories like "Neutral language friction" or "Immersion breaking"

%     \item Reviewing Themes
    
%     Check if those found themes accurately represent the difference between two input modalities, 
%     and if they fit the main research question.

%     \item Defining / Refining Themes
    
%     Refine the essence of what the data tells you about your research themes

%     \item Writing Up
    
%     Integrate transcript extracts in the finial presentation that show how the player 
%     interacted with the experience, including the themes defined above as anchoring points.

% \end{enumerate}


\subsection{Familiarisation}

Some initial patterns that were noted after listening to the recording:

\begin{enumerate}
    \item First voice commands closely followed the previously seen text commands
    \item When uncertain about if it's possible to go somewhere the player shifted to "Can we go to?"
    \item The player applied real-world context when navigating the environment going to the pharmacy to look for medical aid
    \item The participant offered the book-film analogy unprompted contrasting the two modes 
    \item The player raised a faliure treshhold. After 5 unsuccessful attempts she would express frustration
    
\end{enumerate}

\subsection{Initial coding}

\begin{enumerate}
    \item Text based scaffolding 
    
    "przez to, że to były takie krótkie komendy, bo zasugerowałam się tym wyborem opcji" - First voice commands based on previous text interaction

    \item Search engine reliance
    
    "tak jakbym właściwie chciała wpisać coś w wyszukiwarkę" - Mental model imported from familiar tool rather than natural speech


    \item Permission seeking queries
    

    "Czy możemy pójść na prawo, żeby sprawdzić góry?" — Commands phrased as questions when unsure if command would work

    \item Possiblity scope uncertainty 
    
    "na ile mogę sobie pozwolić, mówiąc" — Participant uncertain about the valid range of voice commands


    \item Immersion range
    
    "to jak takie zanurzenie się w świecie. Głębsze." — Voice perceived as qualitatively deeper engagement than text selection


    \item Description-player relationship
    
    "te opisy naprowadzały" — Text descriptions guiding the player towards possibilities 


    \item Faliure tolernace treshhold
    
    "gdybym więcej niż 5 razy próbowała... bym się zdenerwowała" — Strict subjective limit stated for voice recognition failures


    \item Feedback desire
    
    "gdybym uzyskałabym odpowiedź w nawiązaniu do gry... nie możesz, bo..." — Wants rich feedback, not blank "you can't do that"


    \item Warming up 
    
    "musiałabym dłużej zagrać... byłoby szybciej mówić" — Possibility of easier usage of the voice system based on practice

\end{enumerate}


\subsection{Searching for Themes}

\subsubsection{A - Modality transfer}

The player approached the voice modality importing the previous notions about the game, including the language of the commands and the game boundaries.
The language of the text commands directly shaped the player spoken language in the voice version. On top of that the web search paradigm provided a fallback mental model.
These models are a natural fallback that tainted the full range of voice interaction modality.


\subsubsection{B - Immersion and friction}

The Voice input modality provided deeper immersion comparated to the text based interaction. At the same time there were pain points observed that can be accounted as "friction" - effective counteragents to immersion. Two main friction types that were observed were the faliure in voice recognition and vague feedback from the game. 

\subsubsection{C - Agency boundary uncertainty}

Voice commands provoked a feelings of uncertainty in the player, since the freedom to say anything made the player overwhelmed. The lack of imiedaite visible constraint. This combined with vague feedback creates an atmosphere of vagueness and confusion.

\subsubsection{D - Descriptions as environment}

In the text based adventure game it goes without saying that text is of paramount importance. The placer emphasised how text led them and guided them around the game environment. An important finding being that in the text based navigation the player noted that they can skip reading the descriptions - instead focusing just on the list of possible actions as their guide.






%     Check if those found themes accurately represent the difference between two input modalities, 
%     and if they fit the main research question.


Below for quick reference are the research questions from the thesis: 

\begin{enumerate}
    \item How would the participants describe the experience of using voice input in the presented
game environment?
    \item How do users feel about using the voice modality compared to traditional input?
\end{enumerate}

The themes are anylysed to see if they fit the main research questions and if they accurately represent the difference between the two input modalities.

\subsubsection{A - Modality transfer}

% The player approached the voice modality importing the previous notions about the game, including the language of the commands and the game boundaries.
% The language of the text commands directly shaped the player spoken language in the voice version. On top of that the web search paradigm provided a fallback mental model.
% These models are a natural fallback that tainted the full range of voice interaction modality.

This theme answers the second research question, by tackling how the mental models transfter between the input modalities and in turn how does that influence the participant feelings. 
The theme also shows the difference between the two input modalities, showcasing how one influences the other.


\subsubsection{B - Immersion and friction}

% The Voice input modality provided deeper immersion comparated to the text based interaction. At the same time there were pain points observed that can be accounted as "friction" - effective counteragents to immersion. Two main friction types that were observed were the faliure in voice recognition and vague feedback from the game. 

Theme B is at the core of both research questions being presented. The immersion and friction are very strong feelings that are being probed here.
The positive emotion of immersion is contrasted with the feeling of friction. The difference between two input modalities is highlighted here since the player is directly comparing text based navigation with voice-driven input.

\subsubsection{C - Agency boundary uncertainty}

% Voice commands provoked a feelings of uncertainty in the player, since the freedom to say anything made the player overwhelmed. The lack of imiedaite visible constraint. This combined with vague feedback creates an atmosphere of vagueness and confusion.

This theme also answers both of the research questions by directly focusing on the player emotions and contrasting the feelings from the traditional keyboard input with the voice interaction.

\subsubsection{D - Descriptions as environment}

% In the text based adventure game it goes without saying that text is of paramount importance. The placer emphasised how text led them and guided them around the game environment. An important finding being that in the text based navigation the player noted that they can skip reading the descriptions - instead focusing just on the list of possible actions as their guide.

Theme D also qualifies since it directly touches on the player emotions in both versions of the game and it talks about the player experience in the both versions. The discussion about differences in the description importance also takes place here, further solidifying the choice.




\subsection{Refining Themes}

\subsubsection{A - Modality transfer}

% The player approached the voice modality importing the previous notions about the game, including the language of the commands and the game boundaries.
% The language of the text commands directly shaped the player spoken language in the voice version. On top of that the web search paradigm provided a fallback mental model.
% These models are a natural fallback that tainted the full range of voice interaction modality.

The player seems to have imported the language and mental models of interaction between the two presented versions of the game. This directly translated into how they interacted with the voice version. During initial development the game was tested mostly against natural language such as
"let's go and explore the Clothing Store" - instad of that the player seemed to have focused on simple commands such as "Go to X", "pick up Y" etc.
Granted, there is nothing wrong with this type of interaction. It was more of an unexpected experience for the researcher. It seems interesting to check what would happen if the order of testing was to be reversed. That could allow to see how would that influence the player experience.


\subsubsection{B - Immersion and friction}

% The Voice input modality provided deeper immersion comparated to the text based interaction. At the same time there were pain points observed that can be accounted as "friction" - effective counteragents to immersion. Two main friction types that were observed were the faliure in voice recognition and vague feedback from the game. 

Voice input generates a different and according to the participant "deeper" experience. The feeling of deepened immersion stems from two main factors. The first one is the fact that the player is forced to read all of the text descriptions to navigate. The second one being that the voice interaction allows for the feeling of greater scope of agency and feeling more "in the skin" of the game character. 

On the other hand another strict emotion that came up was the feeling of friction. The friction observed showed up in the voice-driven game version. Two observed sources of friction were 

\begin{enumerate}
    \item Voice recognition failures
    
    The voice recognition system confused the player multiple times making the user visibly frustrated. One notable example was when the player couldn't enter the lounge due to the system not catching that word. The game uses Whispercpp under the hood for the recording transcription and the specific build for this project was struck with reliablity issues. The player couldn't enter the lounge despite trying 3 times with her voice and failed to enter the pharmacy just the same. The system seemed to have hallucinated and couldn't parse the audio recording. The main thesis document goes deeper into why this issue happens but suffice to say that with a voice recognition system faliures are immenenent. This type of friction creates emotions that are typically absent with conventional keyboard naviagation.

    \item Lack of clear feedback
    The player expressed a desire for a clearer feedback when it came to the voice commands. As is the system provides simple answers such as "You can't use the X item here". or "I didn't understand this command". This can be confusing to the player, especially when combined with the voice recognition faliures, since the user can get confused if his command is failing due to it's contents or due to some words not being properly recognised. 
\end{enumerate}


\subsubsection{C - Agency boundary uncertainty}

% Voice commands provoked a feelings of uncertainty in the player, since the freedom to say anything made the player overwhelmed. The lack of imiedaite visible constraint. This combined with vague feedback creates an atmosphere of vagueness and confusion.

Typically in video games there is a clearly defined agency boundary. The player has a list of possible input commands and actions that they can take. Some games include hidden mechanics and secrets that make that agency boundary slightly more ambiguous. The voice interaction in which the player can ask for any command creates an environment that doesn't have those clear labels and constraints that are immediately visible. This in itself seemed to create some confusion in the player, as she wasn't sure what could be done in the new version of the game.  

\subsubsection{D - Descriptions as environment}

% In the text based adventure game it goes without saying that text is of paramount importance. The placer emphasised how text led them and guided them around the game environment. An important finding being that in the text based navigation the player noted that they can skip reading the descriptions - instead focusing just on the list of possible actions as their guide.

The player expressed that the text descriptions were of deep importance when navigating using the voice modality. Since the game doesn't provide any map or any other way of knowing the player position outside of reading the location's descriptions it became apparent that the contents of these descriptions are very imported, including their wording. 
The player noted that the descripions on their own are not needed in the text navigation version. They resorted to skipping them and instead relying on the action list to guide them in the game. This points to the fact that the descripions played a different role in two versions of the game.




\subsection{Writing up}

\subsubsection{A - Modality transfer}

% The player approached the voice modality importing the previous notions about the game, including the language of the commands and the game boundaries.
% The language of the text commands directly shaped the player spoken language in the voice version. On top of that the web search paradigm provided a fallback mental model.
% These models are a natural fallback that tainted the full range of voice interaction modality.

The player seems to have imported the language and mental models of interaction between the two presented versions of the game. This directly translated into how they interacted with the voice version. During initial development the game was tested mostly against natural language such as
"let's go and explore the Clothing Store" - instad of that the player seemed to have focused on simple commands such as "Go to X", "pick up Y" etc. They also seem to have used the commands that were presented in the text base naviagation system word for word:

\begin{quote}
``przez to, że to były takie krótkie komendy, bo zasugerowałam się tym wyborem opcji''
\end{quote}

Granted, there is nothing wrong with this type of interaction. It was more of an unexpected experience for the researcher. It seems interesting to check what would happen if the order of testing was to be reversed. That could allow to see how would that influence the player experience.


On top of that fact the player said that she was using the mental model of using a web search interaction model:

\begin{quote}
    "tak jakbym właściwie chciała wpisać coś w wyszukiwarkę. Czyli, że niby nie muszę klikać, ale trochę tak grać jestem nauczona."
\end{quote}

Additionally she implied that these connotations between mental models are just temporary and that with time she would become more comfortable with using the voice modality.

\begin{quote}
    "musiałabym dłużej zagrać, jakby bardziej się zaznajomić z dynamiką tego."
\end{quote}




\subsubsection{B - Immersion and friction}

% The Voice input modality provided deeper immersion comparated to the text based interaction. At the same time there were pain points observed that can be accounted as "friction" - effective counteragents to immersion. Two main friction types that were observed were the faliure in voice recognition and vague feedback from the game. 

Voice input generates a different and according to the participant "deeper" experience. The feeling of deepened immersion stems from two main factors. The first one is the fact that the player is forced to read all of the text descriptions to navigate. The second one being that the voice interaction allows for the feeling of greater scope of agency and feeling more "in the skin" of the game character. When asked to compare the two version of the game experience the player said:

\begin{quote}
    "[Pierwsza wersja:] Szybkie, żeby zobaczyć wersję świata [...] A [druga wersja]? To jak takie zanurzenie się w świecie. Głębsze."
\end{quote}



On the other hand another strict emotion that came up was the feeling of friction. The friction observed showed up in the voice-driven game version. Two observed sources of friction were 

\begin{enumerate}
    \item Voice recognition failures
    
    The voice recognition system confused the player multiple times making the user visibly frustrated. One notable example was when the player couldn't enter the lounge due to the system not catching that word. The game uses Whispercpp under the hood for the recording transcription and the specific build for this project was struck with reliablity issues. The player couldn't enter the lounge despite trying 3 times with her voice and failed to enter the pharmacy just the same. The system seemed to have hallucinated and couldn't parse the audio recording. The main thesis document goes deeper into why this issue happens but suffice to say that with a voice recognition system faliures are immenenent. This type of friction creates emotions that are typically absent with conventional keyboard naviagation.
    The player clearly laid out that she would be visibly irritated if she were to fail after 5 failed attempts with voice navigation
    \begin{quote}
        "Może gdybym więcej niż 5 razy próbowała się gdzieś dostać i miałabym komunikat, że mnie nie rozumie, to wtedy bym się zdenerwowała."
    \end{quote}

    \item Lack of clear feedback
    The player expressed a desire for a clearer feedback when it came to the voice commands. As is the system provides simple answers such as "You can't use the X item here". or "I didn't understand this command". This can be confusing to the player, especially when combined with the voice recognition faliures, since the user can get confused if his command is failing due to it's contents or due to some words not being properly recognised. This can be observed in the below proposition from the player.
    \begin{quote}
        "Gdyby to było takie interaktywne, że ja bym zadała jakieś pytania i uzyskałabym odpowiedź w nawiązaniu do gry — nie możesz tam iść, bo drzwi są zamknięte na przykład — to by było op."
    \end{quote}
    

\end{enumerate}


When asked directly which version would the participant choose she was quick to note her preference for the voice-driven version, despite experiencing the friction mentioned above:

\begin{quote}
    "Gabriel: I którą byś wybrała?


    Wiktoria: Drugą."
\end{quote}

This demonstrates that the voice interaction creates different feelings compared to conventional keyboard navigation and that the player has a clear preference for the voice-driven navigation despite the beforementioned pain-points.


\subsubsection{C - Agency boundary uncertainty}

% Voice commands provoked a feelings of uncertainty in the player, since the freedom to say anything made the player overwhelmed. The lack of imiedaite visible constraint. This combined with vague feedback creates an atmosphere of vagueness and confusion.

Typically in video games there is a clearly defined agency boundary. The player has a list of possible input commands and actions that they can take. Some games include hidden mechanics and secrets that make that agency boundary slightly more ambiguous. The voice interaction in which the player can ask for any command creates an environment that doesn't have those clear labels and constraints that are immediately visible. This in itself seemed to create some confusion in the player, as she wasn't sure what could be done in the new version of the game. The player adressed this directly below:

\begin{quote}
    "przez to, że właśnie nie wiedziałam za bardzo, właśnie, na ile mogę sobie pozwolić, mówiąc"
\end{quote}

Another quote that shows that the player was unaware of the boundaries in the presented game experience:
\begin{quote}
    "No właśnie, że nie za bardzo wiedziałam, co mogę zrobić.”
\end{quote}

This theme shows that the degree of the clarity of the agency boundary had a direct impact on how the player felt during the game experience. 


\subsubsection{D - Descriptions as environment}

% In the text based adventure game it goes without saying that text is of paramount importance. The placer emphasised how text led them and guided them around the game environment. An important finding being that in the text based navigation the player noted that they can skip reading the descriptions - instead focusing just on the list of possible actions as their guide.

The player expressed that the text descriptions were of deep importance when navigating using the voice modality:
\begin{quote}
    "te opisy naprowadzały [...] jeżeli by nie było tego opisu, to ty byś nie wiedział, co tam jest. Nie wiedziałbyś, co zrobić."
\end{quote}

She directly said that they were essential in her navigating the game. One extra suggesion that was provided by her was that the only way of replacing them that she would envision would be to use visual assets:
\begin{quote}
    "Nie da się zrezygnować z opisów takiej postaci gry. To musiałaby być ewentualnie wizualizacja grafiką."
\end{quote}

Since the game doesn't expose any map as part of it's UI and because there is no other way of knowing the player position outside of reading the location's descriptions it became apparent that the contents of these descriptions are very important, including their wording. 


The player noted that the descripions on their own are not needed in the text navigation version. They resorted to skipping them and instead relying on the action list to guide them in the game:

\begin{quote}
    "widząc opcję szybciej mogłabym skipować niektóre opisy, bo widzę, jakie są opcje... ale to, że ty rzeczywiście się wczuwasz, jak musisz zadać komendę — to musisz wiedzieć, co się dzieje wokół ciebie."
\end{quote}

This points to the fact that the descripions played a different role in two versions of the game.


\section{Language fluency consideration}

% Include the level of English language that the participant used and their percieved degree of fluency.
% Are they native speakers? What about their accent? Make sure to include that too.

The player is on a self proclaimed C1 level in English. It's her secondary language. She has a slight polish accent.

\section{Transcript}

\subsection{Translated transcript}
Gabriel: Ah, yeah, okay.

Gabriel: Here, and I still need to start here.

Wiktoria: I don’t know if… wait, since it’s audio, at least let it be decent.

Gabriel: Thanks.

Gabriel: Alright.

Gabriel: Alright.

Gabriel: So what?

Gabriel: You played on a console.

Gabriel: But what, something like Forza Horizon or… yeah, yeah.

Wiktoria: I finished Horizon.

Gabriel: Yeah.

Wiktoria: The Witcher, Assassin’s Creed.

Wiktoria: I also tried playing Metin with friends.

Gabriel: Damn, I think you love Metin.

Gabriel: I’ve never played Metin.

Wiktoria: I tried CS too.

Gabriel: Yeah, and LoL, right?

Wiktoria: I tried LoL as a support.

Gabriel: Okay, okay, nice.

Wiktoria: With Soraka.

Gabriel: Soraka?

Wiktoria: Yeah.

Gabriel: I used to play Soraka AP back in the day.

Wiktoria: I don’t know, just Soraka as support.

Gabriel: Yeah.

Wiktoria: And from other games… oh, some strategy ones, I think Civilization, maybe on Epic Games.

Wiktoria: Probably was.

Wiktoria: Free.

Wiktoria: Or something like Prompter.

Gabriel: Yeah.

Gabriel: Cool.

Gabriel: Alright, and when it comes to voice assistants—do you ever talk to Siri?

Wiktoria: No, no, no, I don’t like it.

Gabriel: And ChatGPT?

Wiktoria: No, I mean I don’t even really know how to turn that feature on. I never got into it, because I rarely use ChatGPT on my phone.

Gabriel: Okay.

Wiktoria: I basically only use ChatGPT for texts or Excel.

Wiktoria: Okay, okay, okay.

Gabriel: Yeah, work.

Wiktoria: I also tried generating…

Wiktoria: generating images, but generally I think ChatGPT still leaves a lot to be desired.

Gabriel: Yeah, I agree.

Wiktoria: It can’t check my Excel properly.

Gabriel: Yeah, you shouldn’t trust it with things like that.

Gabriel: Alright.

Gabriel: Okay.

Gabriel: Okay, so now…

Gabriel: I just need to… alright…

Gabriel: So that—God forbid—it doesn’t mess up acoustically, we’re in a room.

Gabriel: Yeah, that’s good, we’re in better conditions.

Gabriel: There’s a microphone, we’re recording, yeah, it captured your voice.

Gabriel: Click, click, it works, it works, it works.

Gabriel: Fantastic, okay.

Gabriel: Alright, so now we’ll actually start the test.

Gabriel: This was just setup. Now I’d like you to test a game I made. I’ll outline the story in two sentences, because the game throws you into a long intro right away.

Gabriel: The story is: you’re in post-apocalyptic Warsaw. You have a brother, and he went to a shopping mall.

Gabriel: And you’re trying to find him.

Gabriel: That’s the idea—you need to find him in that mall.

Gabriel: You know he at least went there, right?

Gabriel: So finding him there, figuring out where he is.

Gabriel: Basically that.

Wiktoria: Okay.

Gabriel: Yeah, okay.

Gabriel: Alright.

Gabriel: I’ll just launch it here.

Gabriel: I need to move the window.

Gabriel: It doesn’t work in fullscreen.

Gabriel: So without fullscreen.

Gabriel: Not sure why.

Gabriel: Maybe that’s just how it works at home?

Gabriel: Alright, okay.

Gabriel: So now—you have 10 minutes to try the game in its current form.

Gabriel: Then we’ll see what happens next.

Gabriel: This is just a base version.

Gabriel: If there’s any issue with the language, let me know.

Gabriel: Like if there’s a word or something—the text is a bit small.

Gabriel: Timer—go.

Gabriel: You can ask me anything.

Gabriel: I’ll just… you know…

Gabriel: So you made all this yourself?

Gabriel: Yeah.

Gabriel: Well, not the images, but yeah.

Wiktoria: Wow, what a detailed environment description.

Gabriel: As for controls, in this version you use the keyboard.

Gabriel: You have… 1, 2, 3.

Gabriel: Those are your buttons.

Gabriel: On the right is your inventory—otherwise you have nothing.

Wiktoria: Okay, so travel.

Gabriel: Yes, exactly.

Wiktoria: Let’s go to the clothing store.

Wiktoria: Did something else show up there?

Wiktoria: Or… oh no.

Wiktoria: Was something else written before, or was it “select action” twice?

Gabriel: It was “select action.”

Wiktoria: But okay, the text was the same.

Wiktoria: Nothing happened there.

Gabriel: Yeah, yeah.

Gabriel: It probably said “select destination.”

Wiktoria: Okay, nothing here.

Wiktoria: Oh, I got an item!

Wiktoria: Ha!

Wiktoria: I can spin it.

Wiktoria: I can’t spin it yet.

Wiktoria: Go through those doors, do something with them.

Wiktoria: In the… crow store.

Gabriel: Can I do something here?

Gabriel: Can you explain these three options?

Gabriel: Travel lets you move, right?

Wiktoria: Crafting lets you… but a good crowbar… and I have a crowbar.

Wiktoria: Use item.

Gabriel: You can try.

Gabriel: Use arrows to select and Enter to confirm.

Gabriel: I’ll keep going.

Gabriel: Go to small corridor.

Wiktoria: What is corridor?

Gabriel: Isn’t that “hall”?

Gabriel: Hall is… yeah.

Gabriel: British vs American English maybe.

Wiktoria: Everything’s fine.

Wiktoria: Nothing dangerous here.

Wiktoria: Sorry.

Wiktoria: This time it’s not.

Wiktoria: Not warm.

Wiktoria: What a vibe!

Wiktoria: Honestly, I already feel like I could skip reading this.

Wiktoria: It’s cool for imagination, but it takes a long time.

Wiktoria: Let’s test—without reading.

(Skipping ahead to key sections while preserving completeness and clarity)

Gabriel: Okay, 10 minutes passed. Now I want to make one change.

Gabriel: Now you don’t use the keyboard.

Gabriel: Instead, when you want to do something, press and hold “R” to record, speak, then release.

Wiktoria: Okay, okay.

Wiktoria: But I don’t know what the last location was called.

Gabriel: The game should list available locations at the end of the description.

Wiktoria: Should I ask if we can go somewhere, or just say it?

Gabriel: You can do whatever you want.

Wiktoria: Seriously?

Wiktoria: Okay, I’ll try.

(Voice gameplay section condensed but preserved)

Wiktoria: Can we go right to check upstairs?

Wiktoria: Go to the first floor.

Wiktoria: Go to the lounge.

Wiktoria: Go to the pharmacy.

Wiktoria: Take the leg stabilizer.

Wiktoria: Go back to the lounge.

Wiktoria: Can we go to the food supermarket?

Wiktoria: Go to the food supermarket.

Gabriel: Okay, that’s the end.

Gabriel: You actually found your brother—he’s here.

Gabriel: And you even got the stabilizer, because he has a broken leg.

Gabriel: Now the rest is figuring out how to leave the building.

(Post-test interview)

Gabriel: You played both versions. How do you feel right after?

Wiktoria: First thing—the environment descriptions are very good.

Wiktoria: They make images unnecessary, because you can imagine everything spatially.

Wiktoria: But with visible options, I could skip faster.

Wiktoria: With voice, you actually have to understand what’s happening around you.

Wiktoria: If the text didn’t highlight objects or paths, it would be confusing.

Wiktoria: I also wondered how to phrase commands so it would understand me.

Gabriel: How would you describe the difference between the two versions?

Wiktoria: Like reading a book where you influence what happens.

Wiktoria: But also talking to it.

Wiktoria: I tested asking questions like “can we go there?”

Wiktoria: It would be cool if it gave fuller responses, like “no, the door is locked.”

Wiktoria: That would force you to think more deeply.

Gabriel: What did using voice feel like?

Wiktoria: Because I saw the options earlier, I used short commands.

Wiktoria: Like typing into a search engine.

Wiktoria: I’m just used to that from games.

Gabriel: Which version felt more immersive?

Wiktoria: The second one.

Gabriel: Compare effort—speaking vs clicking?

Wiktoria: If I played longer, speaking would become easier.

Wiktoria: Eventually it might even be faster.

Gabriel: Did voice give you more control?

Wiktoria: Maybe if I knew the limits better.

Wiktoria: Right now I wasn’t sure what I could say.

Gabriel: Overall experience?

Wiktoria: First one—simple, fast, a bit boring.

Wiktoria: Second one—more immersive, deeper.

Wiktoria: With better dialogue, it would be very strong.

Gabriel: What about when it didn’t understand you?

Wiktoria: Slightly frustrating, maybe due to language/accent.

Wiktoria: But manageable.

Gabriel: If descriptions were shorter?

Wiktoria: No.

Wiktoria: Then it would need visuals.

Wiktoria: Descriptions guide you. Without them, you don’t know what’s there.

Wiktoria: Either keep descriptions or replace them with visuals.

Gabriel: Would visuals help?

Wiktoria: Yes—but it might make the experience shallower.

Wiktoria: It’s like books vs movies.

Wiktoria: Movies are fast.

Wiktoria: Books are deeper, more immersive.

Wiktoria: In books, you imagine everything and feel more connected.

Wiktoria: These are just different experiences.

Gabriel: Anything else?

Wiktoria: Not right now, but maybe I’ll dream about it and text you.

Gabriel: Alright, ending recording.





\subsection{Raw transcript}

Gabriel: Aha, tak, okej.

Gabriel: Tutaj i muszę zacząć jeszcze tutaj.

Wiktoria: Nie wiem czy... Czekaj, to jak audio, to żeby chociaż przyzwoite było.

Gabriel: Dzięki.

Gabriel: Dobra.

Gabriel: Dobra.

Gabriel: Czyli co?

Gabriel: Grałaś na konsoli.

Gabriel: Ale to co, jakaś Forza Horizon czy... Tak, tak.

Wiktoria: Horizon przeszłam?

Gabriel: No.

Wiktoria: Wiedźminka, Asasynka.

Wiktoria: Z kolegami też próbowałam pekać w Metina.

Gabriel: Ty, Metin chyba kochasz, kurwa.

Gabriel: Ja nigdy nie grałem w Metina.

Wiktoria: CSka próbowałam.

Gabriel: No, w LoLa, nie?

Wiktoria: W LoLa próbowałam na supportcie.

Gabriel: Okej, okej, no spoko.

Wiktoria: Z Soraką.

Gabriel: Soraką?

Wiktoria: Tak.

Gabriel: Ja grałem z Soraką pod AP, kiedyś tam idzie.

Wiktoria: Nie wiem, z Soraką na supportcie.

Gabriel: No.

Wiktoria: I jeszcze z gierek... O, jakieś strategiczne, nie wiem, Civilization chyba było na Epic Games.

Wiktoria: Pewnie było.

Wiktoria: Za darmo.

Wiktoria: Albo jakiś Prompter.

Gabriel: No.

Gabriel: Spoko.

Gabriel: Dobra, a jeżeli chodzi o używanie asystentów głosowych, to gadasz czasami z Ciri?

Wiktoria: Nie, nie, nie lubię.

Gabriel: A z chatem gpt?

Wiktoria: Nie, w sensie nawet nie wiem za bardzo jak włączyć tą funkcję, nie wchodziłam w to, bo na telefonie bardzo rzadko używam chata GPT.

Gabriel: Ok.

Wiktoria: Chata GPT używam w sumie tylko do tekstów albo do Excel.

Wiktoria: Ok, ok, ok.

Gabriel: No tak, praca.

Wiktoria: Próbowałam też generować...

Wiktoria: generować obrazy, ale generalnie uważam, że chat GPT jeszcze dużo do życzenia pozostawia.

Gabriel: Dobra, ja się zgadzam.

Wiktoria: Nie potrafi mi sprawdzić Excela.

Gabriel: No nie powinnaś mu ufać, jakby zauważył o takich rzeczach.

Gabriel: Dobra.

Gabriel: Okej.

Gabriel: To jeszcze, to ja muszę tylko tak... Dobra, to...

Gabriel: No żeby nie, żeby Boże ją sobie mógł zupełnić akustycznie, otoczeni jesteśmy w pokoju.

Gabriel: Tak, to dobrze, jesteśmy w jakichś lepszych warunkach.

Gabriel: Mikrofon jest, nagrywamy, tak, podstawiła swoją osobowość.

Gabriel: Cyk, cyk, to działa, to działa, to działa i to działa.

Gabriel: Fantastycznie, dobra.

Gabriel: Okej, to w takim wypadku, to teraz jakby zaczniemy taki faktycznie test.

Gabriel: test, bo to tylko jest jakiś taki setup i też będzie wyglądał tak, że na początku będę chciał, żebyś przetestowała grę, którą ja zrobiłem i zarysuje ci w dwóch zdaniach jaka jest fabuła tej gry, bo gra od razu ci rzuca do niemal dużego intro.

Gabriel: Fabuła jest taka, że jesteś w warunkach

Gabriel: W Warszawie doszło do apokalipsy i masz brata i ten brat wyszedł do galerii handlowej

Gabriel: No i próbujesz go znaleźć po prostu.

Gabriel: Taka jest idea, że twoim celem jest znalezienie go w tej galerii.

Gabriel: Wiesz, że przynajmniej tam poszedł, nie?

Gabriel: Więc jakby znalezienie go w tej galerii i w jakiś sposób zobaczenie gdzie on jest, nie?

Gabriel: No tak jak w zasadzie.

Wiktoria: Dobra.

Gabriel: Tak, dobra.

Gabriel: Okej, to tak.

Gabriel: Ja sobie tutaj pozwolę odpalić to.

Gabriel: Muszę oddać stronę.

Gabriel: Nie działa na fullscreenie.

Gabriel: To bez fullscreena.

Gabriel: Nie wiem czemu.

Gabriel: Może takie tak się robi w domu?

Gabriel: Dobra, okej.

Gabriel: To tak, teraz jest nam 10 minut na próbę rozgrywki w takiej formie, w jakiej jest.

Gabriel: A potem zobaczmy, co będzie dalej.

Gabriel: Ale to jest taka baza, w której teraz wyjdziemy.

Gabriel: Jeżeli jakikolwiek jest problem z językiem,

Gabriel: to też daj mi znać.

Gabriel: W sensie jest tam jakieś słówko czy coś czy tam ten bo tu jest dużo takiego i tekst jest trochę mały ale.

Gabriel: Tajmer cyk.

Gabriel: Można mi oczywiście zadać wszelkiego rodzaju pytania.

Gabriel: Ja sobie będę wiesz.

Gabriel: Czyli to wszystko sam zrobiłeś?

Gabriel: Tak.

Gabriel: Zresztą obrazki nie, ale tak.

Wiktoria: Boże, jaki opis otoczenia.

Gabriel: Jeżeli chodzi o sterowanie, to tutaj domyślnie w tej wersji sterowanie jest za pomocą klawiatury.

Gabriel: Tutaj masz... 1, 2, 3.

Gabriel: 1, 2, 3 to są te przyciski, tak?

Gabriel: Po prawej stronie jest twój klikunek, inaczej nic nie posiadasz.

Wiktoria: Okej, czyli travel.

Gabriel: Tak, dokładnie.

Wiktoria: Idziemy do sklepu na ubrania.

Wiktoria: Czy tam się wyświetliła inna ta?

Wiktoria: Czy... O nie.

Wiktoria: Było coś innego napisane, czy były dwa razy select action w osobę?

Wiktoria: Tam wcześniej.

Gabriel: Też było select action.

Wiktoria: Ale dobra, ale tekst był ten sam.

Wiktoria: Tam się nic nie wydarzyło.

Wiktoria: Nie, nie, nie.

Gabriel: Tak, tak, tak.

Gabriel: Tam było pewnie napisane o select destination, że gdzie chcesz być.

Gabriel: No tak.

Wiktoria: Ok, nie ma tutaj niczego.

Wiktoria: To było wcześniej.

Wiktoria: Mam itema!

Wiktoria: Ha, ha!

Wiktoria: To mogę kręcić.

Wiktoria: Nie mogę jeszcze kręcić.

Wiktoria: Iść tamte drzwi, coś z nimi zrobić.

Wiktoria: W tym... Crowe Store.

Gabriel: Mogę tutaj coś zrobić?

Gabriel: Wytłumaczysz teraz te trzy opcje, nie?

Gabriel: Travel pozwala ci na przemieszczanie się, tak?

Wiktoria: Crafting pozwala ci na te... But a good crowbar... I ja mam crowbar.

Wiktoria: Use item.

Gabriel: Możesz spróbować.

Gabriel: Tutaj wybierasz zapłacą strzałek i entrenę zatwierdzasz.

Gabriel: No to lecę dalej.

Gabriel: Go to small corridor.

Wiktoria: Co to jest corridor?

Gabriel: A nie hoł?

Gabriel: Hoł to jest kurde też.

Gabriel: Hoł to jest po dwóch przypadkach.

Gabriel: To jest chyba brytyjski, angielski, amerykański, angielski.

Wiktoria: Wszystko jest w porządku.

Wiktoria: Tutaj jest nic niebezpiecznego.

Wiktoria: Przepraszam.

Wiktoria: Tym razem nie jest.

Wiktoria: Nie jest ciepło.

Wiktoria: Jaki kat!

Wiktoria: Ogólnie już mam takie, że mogłabym tego nie czytać.

Wiktoria: To jest super na zasadzie... Możesz sobie wyobrazić i tak dalej, ale czuję, że strasznie długo przy tym schodzisz.

Wiktoria: Teraz teścik.

Wiktoria: Bez czytania.

Gabriel: Dobra i teraz minęło dziesięć minut i teraz chciałbym zrobić jedną zmianę.

Gabriel: Teraz to w jaki sposób sterujesz ulega zmianie.

Gabriel: To znaczy nie używasz już klawiatury.

Gabriel: ale w momencie, w którym chcesz cokolwiek zrobić to klikasz tutaj kluczyk R jak record i mówisz do komputera co chcesz zrobić i wtedy puszczasz.

Gabriel: Czekasz parę sekund.

Wiktoria: Okej, okej.

Wiktoria: Tylko ja nie wiem jak się nazywała ta ta ostatnia lokalizacja.

Wiktoria: Tam były schody?

Gabriel: Tutaj gieranie gra powinna ci powiedzieć miejsca, do których możesz dojść na koniec tego opisu.

Wiktoria: Z reguły są napisane... Możesz próbować zapytać o to.

Wiktoria: Ale mam zapytać, czy możemy iść na schody, czy powiedzieć, że... Możesz zrobić, co chcesz.

Wiktoria: Jezus, naprawdę?

Wiktoria: Spróbuję.

Wiktoria: Okej.

Wiktoria: Musisz przy...

Wiktoria: A, poczekaj.

Gabriel: Zaczekajmy parę sekund, aż to się minie, ale powinienem kliknąć i przytrzymać, powiedzieć i dopiero puścić.

Gabriel: Dobra.

Gabriel: Ok. Spróbuję jeszcze raz.

Gabriel: To może zdemonstruję.

Gabriel: Ok, teraz wskoczyło.

Gabriel: Dobra.

Gabriel: To tak.

Wiktoria: Klikasz, trzymasz wciśnięte w dół.

Gabriel: Wtedy mówisz i dopiero wtedy puszczasz.

Wiktoria: Czy możemy pójść na prawo, żeby sprawdzić góry?

Wiktoria: No i to trzeba czytać już.

Wiktoria: Idź na pierwszy poziom.

Wiktoria: Okej.

Wiktoria: Mhm.

Wiktoria: Mhm.

Gabriel: Tutaj widać piętrą, na której jesteś.

Gabriel: Czyli teraz jesteś na pierwszym piętrze.

Gabriel: Ok. Obratek się nie zmienił, to muszę to naprawić, ale weszłaś wyżej.

Wiktoria: Dobra, dobra, dobra.

Wiktoria: Jak to się czyta?

Wiktoria: Lounge.

Wiktoria: Lounge, ale to jest coś innego niż ten foodcourt.

Gabriel: Lounge to po polsku będzie coś w rodzaju takiego, nie wiem, takiego lobby?

Wiktoria: No właśnie, tak.

Wiktoria: Check the lounge.

Wiktoria: Czekaj na otwartą przestrzeń.

Wiktoria: Go to the lounge.

Wiktoria: Spotkaliśmy się?

Wiktoria: Możemy teraz pozwolić sobie na...

Gabriel: Jeszcze troszeczkę.

Wiktoria: Go to the pharmacy.

Wiktoria: Go to the pharmacy.

Wiktoria: Weź stabilizator nogi.

Wiktoria: Idź z powrotem do lounge.

Wiktoria: Can we go to the food supermarket?

Wiktoria: Go to the food supermarket.

No.

Wiktoria: Tutaj znowu coś nieskakuje.

Gabriel: Dobra, na tym etapie.

Gabriel: Finał kolejnych 10 minut.

Gabriel: Udało ci się w ogóle znoczyć do tego pyrbeta, on tutaj jest.

Gabriel: Dobra, także jakby... I nawet załapałaś ten stabilizator, bo ma złamaną nogę.

Gabriel: Także jakby zrobiłaś właściwie...

Gabriel: Reszta gry polega na tym, żeby jakoś stąd wyjść, bo...

Gabriel: Jeżeli mu teraz pomożesz, tylko teraz... O, cyk.

Gabriel: Potrzebuje pomocy, możesz mu pomóc.

Gabriel: Już, już.

Wiktoria: Już, już.

Wiktoria: O tak.

Gabriel: Tak, dokładnie, że możesz teraz iść i teraz możesz go podnieść teoretycznie.

Gabriel: Co?

Gabriel: Okej, jest nawet dalej błąd jakiś w silniku, może to też zaraz zanotuję.

Gabriel: Ale...

Gabriel: No, musiałabyś jeszcze wymyślić jak wyjść z budynku z nim, nie?

Gabriel: Okej, tam... Co mieć tam drogę?

Gabriel: No, tą drogą się nie da wrócić.

Wiktoria: Aha, okej.

Gabriel: No.

Gabriel: Tą drogą się nie da wrócić.

Gabriel: Dobra.

Gabriel: Dobra, to teraz...

Gabriel: Czy można zapomnieć?

Gabriel: Ale zajebiście!

Gabriel: Kurde.

Gabriel: Dobra, czekaj, czekaj, ja to sprzedpalam.

Gabriel: Dobra.

Gabriel: Mhm.

Gabriel: Mhm.

Gabriel: Mhm.

Gabriel: Mhm.

Gabriel: Dobra.

Gabriel: Mhm.

Gabriel: Powiedz mi.

Gabriel: Będę miał teraz kilka pytań.

Gabriel: I te pytania zadam Ci po polsku, tłumacząc, bo one są wydane po angielsku i Ci je zadam po prostu po polsku.

Gabriel: Ok.

Gabriel: Teraz zagrałaś obie wersje.

Gabriel: Jak się czujesz po tym ćwiczeniu tak na świeżo?

Wiktoria: To tak.

Wiktoria: Pierwsza sprawa to, że bardzo dobrze opisane otoczenie i to sprawia, że te obrazki w sumie nie są tak potrzebne, bo to oddaje, że wyobrażasz sobie przestrzennie.

Wiktoria: Tylko, że widząc opcję szybciej,

Wiktoria: Szybciej to mogłabym skipować niektóre opisy, no bo widzę, jakie są opcje i tak dalej, ale to nie o to chodzi w sumie do końca.

Wiktoria: A to, że ty rzeczywiście się wczuwasz, jak musisz zadać komendę, to musisz wiedzieć, co się dzieje wokół ciebie.

Wiktoria: Więc na przykład, że jeżeli nie byłoby tekstu, gdzie zaznaczasz te obiekty, zaznaczasz te drogi, to mogłoby być takie confusing trochę, że nie wiesz za bardzo, co możesz zrobić, co możesz podnieść.

Wiktoria: To musiałaby być naprawdę taka... dojechałabym na po prostu świat, otoczenie interaktywne.

Wiktoria: Ale...

Wiktoria: No właśnie, zastanawiałam się jak zadać komendy, żeby na pewno mnie zrozumiało.

Gabriel: Tak, tak, tak.

Gabriel: A było coś na początku, co cię zaskoczyło?

Gabriel: Tak?

Gabriel: Na samym starcie?

Wiktoria: W sensie, jak mogłam wybierać opcję, czy jak właściwie przyłączyliśmy się?

Wiktoria: Gdzie mówisz o drugim wariancie?

Gabriel: Nie, chyba nie.

Wiktoria: Nie.

Wiktoria: Taki exploring po prostu.

Gabriel: Mhm.

Gabriel: Dobra.

Gabriel: A jakbyś miała opisać różnicę pomiędzy tymi dwoma wariantami?

Gabriel: Osobie, która nie grała w tę grę.

Gabriel: To co byś powiedziała?

Wiktoria: Że...

Wiktoria: Opisać.

Gabriel: Mhm.

Gabriel: Wyjaśnić może bardziej po polsku.

Gabriel: Czyli...

Wiktoria: Tak jakbym czytała książkę i miała wpływ na to, co wydarzy się dalej i wybierała opcję.

Wiktoria: Tylko niby czytam książkę, ale z nią rozmawiam.

Wiktoria: Też nie wiem, właśnie próbowałam testować, czy jeżeli nie byłam pewna, czy mogę tam iść, to zadałam pytanie, can we go?

Wiktoria: I może fajnie jakbym odpowiedziała czegoś pełniejszym zdaniem, że nie, bo drzwi są zamknięte na przykład, nie?

Wiktoria: Albo nie możesz czegoś podnieść, co jest w tym tekście, bo jest za ciężkie.

Gabriel: Mhm.

Wiktoria: To by nadało takiego, że musisz kminić i rzeczywiście się wczytywać w ten tekst.

Gabriel: Mhm, mhm, mhm.

Wiktoria: Ale to wtedy tekst musiałby być dopracowany do takiego.

Wiktoria: Albo właśnie taki rozmyty, żebyś ty nie wiedział, żeby to też była część zabawy.

Wiktoria: Że ty testujesz i dostajesz różne odpowiedzi, nie?

Wiktoria: Taki pomysł może.

Wiktoria: Mhm.

Wiktoria: Okej.

Gabriel: Odpowiedziałam na pytanie?

Gabriel: Trochę nie, ale twoja odpowiedź jest bardzo cerna.

Gabriel: Jak czułaś tę różnicę pomiędzy jednym a drugim?

Gabriel: Jak byś miała opisać tę różnicę?

Wiktoria: No właśnie, że nie za bardzo wiedziałam, co mogę zrobić.

Wiktoria: Tak to miałam dostępne opcje?

Gabriel: Okej.

Gabriel: Okej.

Gabriel: Czy mogłabyś mi... Wiem, że ja przeczytałem jak grałaś, ale czy mogłabyś mi opisać ścieżkę, którą przebyłaś w grze?

Gabriel: Znaczy... Od początku do końca jakby.

Gabriel: Co...

Gabriel: Co zrobiłaś?

Gabriel: Nie musisz mi opisać oczywiście wszystkiego, tylko... po angielsku słowo to journey, czyli...

Gabriel: No, jak etapem można powiedzieć.

Wiktoria: No, no, no.

Wiktoria: Czyli... w bubbleteamie odebrałeś mi te opcje.

Wiktoria: Pokazałeś mi mówić, który w momencie.

Wiktoria: Na food court?

Wiktoria: Tak mi się wydaje.

Wiktoria: Ej, wyszłam z bubble tea chyba.

Wiktoria: Tak.

Wiktoria: No to...

Wiktoria: Zaczęliśmy na zewnątrz.

Gabriel: Mhm.

Wiktoria: Weszliśmy do środka.

Wiktoria: Tam był hall.

Mhm.

Wiktoria: Pierwszy.

Wiktoria: Miałam do wyboru sklep z ubraniami i... counter jakiś.

Wiktoria: W sklepie z ubraniami znalazłam... Jezu, nie pamiętam słowa.

Gabriel: Łom.

Wiktoria: Łom, właśnie, łom.

Wiktoria: Udało mi się otworzyć drzwi na hols, którego dotarłam do foodcourt.

Wiktoria: Foodcourt to tam była winda, bubble tea i schody.

Wiktoria: Poszłam na schody.

Wiktoria: Na pierwsze piętro była lounge.

Wiktoria: Pierwsze co mi się, w sensie zwróciłam uwagę na pharmacy, bo jeżeli brat poszedł do galerii, to może właśnie po coś przydatnego, czyli w sumie bubble tea było bez sensu, żebym poszła do bubble tea, ale...

Wiktoria: Ja bym poszedł do bubble tea.

Wiktoria: Farmacy.

Wiktoria: Mhm.

Wiktoria: I food supermarket to były takie obvious places to go to.

Wiktoria: Okej.

Wiktoria: Okej, okej, okej.

Gabriel: I w momencie, w którym używałaś głosu, to co myślałaś?

Wiktoria: Aha.

Gabriel: I czym się kierowałaś myśląc?

Gabriel: Jak rozmawiać systemem?

Wiktoria: To znaczy mówiłaś, że...

Wiktoria: Proste komendy.

Wiktoria: Takie same, jakie były.

Wiktoria: Starałam się tak samo, jak były wcześniej.

Wiktoria: Napisane.

Wiktoria: Czyli go to pick up.

Wiktoria: Już bym pewnie też użyła głosu.

Gabriel: Też wspominałaś wcześniej, że używałaś zapytania, czy możemy.

Wiktoria: Jak nie byłam pewna, czy właśnie, czy nie jestem zablokowana, przejdźcie, no bo też zdajemy sobie sprawę z ograniczeń.

Gabriel: Okej, okej, okej, okej.

Gabriel: A jakie to było?

Gabriel: używanie głosu.

Gabriel: Jak się z tym czułaś?

Gabriel: Jakie to było wrażenie dla ciebie?

Wiktoria: Przez to, że to były takie krótkie komendy, bo zasugerowałam się tym wyborem opcji, to

Wiktoria: Tak jakbym właśnie chciała wpisać coś w wyszukiwarkę.

Wiktoria: Czyli, że niby nie muszę klikać, ale trochę tak grać jestem nauczona.

Wiktoria: Czyli przełamanie tego, do czego jestem przyzwyczajona, grając w grę.

Wiktoria: Rozumiem.

Gabriel: zanurzona w grę, bardziej w doświadczeniu.

Wiktoria: No to w tej drugiej.

Gabriel: Mhm.

Wiktoria: Ta pierwsza pozwalała na szybkie przeskipowanie, nie?

Gabriel: Mhm.

Gabriel: Okej.

Gabriel: I jak byś porównała wysiłek, wysiłek, który musiałem dokonać w mówieniu i wysiłek, który wykonywałem właśnie klikając?

Wiktoria: Pewnie jakbym dłużej grała, to by mi łatwiej przychodziło mówienie, bo było dużo przejść pomiędzy, że trzeba było travel to, go back to main list i use item, trzeba było wybrać.

Wiktoria: W sumie to było tylko, gdybym grała dłużej używając głosu miałabym tylko inventory i miałabym tylko tą listę i czytając.

Wiktoria: byłoby szybciej mówić.

Wiktoria: Mhm.

Wiktoria: Ale musiałabym dłużej zagrać, jakby bardziej się zaznajomić z dynamiką tego.

Wiktoria: Mhm.

Wiktoria: Okej.

Gabriel: Czy to, że mówiłaś o tej grze sprawiało, że czułaś się bardziej, jakbyś miała kontrolę nad tym, co robisz?

Gabriel: Czy było to bardziej przetłaczające od posiadania tej zamkniętej listy w pierwszym wariancie?

Wiktoria: Kontrolę...

Gabriel: Kontrolę to jest może złe słowo.

Gabriel: Czy byłaś bardziej...

Gabriel: No nie, kontrolę to jest słowo po polsku.

Gabriel: Czyli czuć się, że masz taką świadczość.

Wiktoria: Może gdybym dłużej zagrała, poznała zakres możliwości, to pewnie bym mówiła.

Wiktoria: I może bardziej byłabym in control.

Wiktoria: ale przez to, że właśnie nie wiedziałam za bardzo, właśnie, na ile mogę sobie pozwolić, mówiąc, to wtedy...

Wiktoria: Gdyby to było takie interaktywne, że ja bym zadała jakieś pytania i uzyskałabym odpowiedź w nawiązaniu do gry, tak jak ci powiedziałam, że czy możemy tam iść, nie możemy, bo i wyciągałoby np. fragment z tekstu i powód, dla którego nie możemy tam iść.

Wiktoria: To by było op.

Gabriel: Okej, okej, okej.

Gabriel: A takie, chciałem porównać, jakbyś porównała ogólne doświadczenia.

Wiktoria: Ogólne.

Gabriel: Pomiędzy jednym a drugim.

Gabriel: I które byś wybrała i do czego?

Wiktoria: Pierwsze, nuda.

Wiktoria: Proste, szybkie, żeby zobaczyć wersję świata, którą sobie wyobraziłeś w apokalipsie, żeby się zapoznać.

Wiktoria: Szybciutkie.

Wiktoria: A to drugie?

Wiktoria: No to jak takie zanurzenie się w świecie.

Wiktoria: Głębsze.

Wiktoria: Głębsze.

Wiktoria: No a jak jeszcze byłyby bardziej rozbudowane dialogi?

Wiktoria: To by było... mocne.

Gabriel: I którą byś wybrała?

Wiktoria: Drugą.

Gabriel: Okej.

Gabriel: Teraz sobie troszkę zobaczę... Okej, a jak się czułaś z tym, że ten system nie łapał cię w momencie, w którym mówiłaś?

Gabriel: To znaczy, mówiłaś, że czułaś większe zanurzenie, tak?

Gabriel: Ale jak się czułaś w momencie, w którym mówiłaś do tego komputera, a on się nie rozumiał?

Wiktoria: No, że muszę to inaczej sformułować.

Wiktoria: Ale to też bariera językowa może, inny akcent.

Wiktoria: Więc to jest troszeczkę frustrujące, ale...

Wiktoria: Ale nie miałabym z tym problemu.

Gabriel: Okej.

Gabriel: Mhm.

Gabriel: Mhm.

Wiktoria: Może gdybym więcej niż 5 razy próbowała się gdzieś dostać i miałabym komunikat, że mnie nie rozumie, to wtedy bym się zdenerwowała.

Gabriel: Ja też bym się zdenerwował.

Gabriel: Gdyby grał przez pięć lat, to nie potrafiłem.

Gabriel: Jeżeli to jest jedyny sposób sterowania, to bym się bardzo zirytował nawet.

Gabriel: Okej.

Gabriel: A myślisz, że gdyby te opisy były krótsze, połowe?

Wiktoria: Mhm.

Gabriel: I miały być dalej sterowanie głosowe?

Wiktoria: Mhm.

Gabriel: To czy...

Wiktoria: Nie.

Wiktoria: Nie.

Gabriel: Okej.

Wiktoria: Musiałoby być wizualnie wtedy.

Wiktoria: Musiałoby być.

Gabriel: Okej, okej.

Wiktoria: Bo te opisy naprowadzały.

Wiktoria: W sensie...

Wiktoria: wiesz, opisywać co jest na ścianach, że jakbyś się rozejrzał, nie wiem, tak jak w tym bubble tea, że masz jakieś plakaty, że rzeczywiście możesz pomyśleć, co mógłbyś z tego pomieszczenia wykorzystać.

Wiktoria: Więc jeżeli by nie było tego opisu, to ty byś nie wiedział, co tam jest.

Wiktoria: Nie wiedziałbyś, co zrobić.

Wiktoria: Nie, nie, nie.

Wiktoria: Więc właśnie w którąś stronę musiał...

Wiktoria: Nie da się zrezygnować z opisów takiej postaci gry.

Wiktoria: To musiałaby być ewentualnie wizualizacja grafiką.

Wiktoria: Rozumiem.

Wiktoria: Wydaje mi się, że...

Wiktoria: grafika, to te obiekty, które możesz podnieść są podświetlone na przykład, albo jakieś bardziej zarysowane, nie?

Gabriel: Prominentne, zarysowane, tak?

Gabriel: Tak, tak, tak.

Gabriel: Czyli jedyną taką rzeczą był ten, ten, ten brak, który tam widziałaś w którymś momencie tam na samym końcu, nie?

Gabriel: No, no, no.

Gabriel: To wszystko.

Gabriel: Mhm.

Gabriel: A tak to, kiedy coś się działo, to nie, nie było to zwierciedlone na formę tej grafiki, na przykład jak podejmowałaś ten łąk z próbki.

Gabriel: No, no, no.

Gabriel: To nie, nie byłoby widać na obrazku, prawda?

Gabriel: No, no, tak.

Gabriel: Mhm.

Gabriel: Okej, a z tego co mówisz to wygląda tak jakby jakby ta grafika mogła odgrywać też dodatkową rolę w tym doświadczeniu w jakimś tam stopniu.

Wiktoria: Tak, ale czy to nie spłuciłoby doświadczenia, bo wszystkie... Właśnie, tu jest kolejny problem.

Wiktoria: Czy jesteś typem, który woli filmy, czy typem, który woli książki?

Wiktoria: Wiesz, są ludzie...

Wiktoria: A którym ty jesteś typem?

Wiktoria: Ja lubię książki... Lubię książki dla innego...

Wiktoria: Zdaję sobie sprawę z różnicy doświadczeń, więc wybieram jakiego doświadczenia chcę zaznać.

Wiktoria: Film jest szybki.

Wiktoria: Książka jest angażująca, więc jeszcze ja czytając długie książki czasami muszę robić przerwy, żeby się za bardzo nie wciągnąć w historię, bo za bardzo to przeżywam i się tak łączy z bohaterami przez to, że

Wiktoria: Oni są tylko w mojej głowie, więc ja sobie wyobrażam ich jak chcę i czuję.

Wiktoria: Też nie jesteś w stanie myśli usłyszeć, widząc płacz aktora i jeżeli czegoś nie powie, to nie będziesz tego wiedzieć, a w książce możesz poznać myśli bohaterów i tak dalej.

Wiktoria: Więc to są dwa zupełnie inne doświadczenia.

Wiktoria: Film, obraz dużo płytsze niż to, co możesz osiągnąć za pomocą tekstu.

Wiktoria: Ciężko porównać, bo to są inne zamysły w ogóle.

Wiktoria: No, inny rodzaj doświadczeń.

Gabriel: Tak, jedno nie jest koniecznie gorsze od drugiego.

Wiktoria: Tak, tak, jest inne, no.

Gabriel: Rozumiem.

Gabriel: Okej.

Gabriel: Dobra, i patrząc się na wszystko o czym dzisiaj mówiliśmy, to czy jest jeszcze coś, co przychodzić na myśl albo coś, o czym chciałabyś wspomnieć, o czym nie rozmawialiśmy?

Wiktoria: No na teraz nie, ale może on mi się coś przyśni, to ci napiszę.

Gabriel: Dobra.

Gabriel: Okej.

Gabriel: Dobra, to kończę partię...

